\chapter*{Lista de Siglas e Termos Técnicos}
\addcontentsline{toc}{table}{Lista de Siglas e Termos Técnicos}
\lhead{LISTA DE SIGLAS E TERMOS TÉCNICOS}

\begingroup
\scriptsize
\renewcommand{\arraystretch}{1.22}
\setlength{\tabcolsep}{3pt}
\begin{longtable}{@{}>{\raggedright\arraybackslash}p{2.35cm}>{\raggedright\arraybackslash}p{5.05cm}>{\raggedright\arraybackslash}p{7.15cm}@{}}
\toprule
\textbf{Termo/Sigla} & \textbf{Significado} & \textbf{Interpretação no trabalho} \\
\midrule
\endfirsthead
\toprule
\textbf{Termo/Sigla} & \textbf{Significado} & \textbf{Interpretação no trabalho} \\
\midrule
\endhead
\endfoot
AIPE & \textit{Adaptive Inventory Policy Engine} & \textit{Framework} metodológico que recomenda políticas de reposição por série loja-produto. \\
PSE & \textit{Policy Selection Engine} & Meta-modelo supervisionado que aprende a relação entre características da série e política recomendada. \\
Política de reposição & Regra ou algoritmo de decisão de estoque & Define quando emitir pedido e em qual quantidade. \\
Série loja-produto & Histórico de vendas de um produto específico em uma loja específica & Unidade básica de análise da dissertação. \\
POD & Perfil Operacional de Demanda & Agrupamento ou caracterização do padrão operacional da série. \\
CTI & Custo Total de Inventário & Métrica de custo usada na avaliação das políticas; deve ser minimizada. \\
NS & Nível de Serviço & Proporção da demanda atendida; deve ser maximizada. \\
TR & Taxa de Ruptura & Frequência ou proporção de falta de estoque; deve ser minimizada. \\
BE & \textit{Bullwhip Effect} & Amplificação da variabilidade dos pedidos na cadeia. \\
FP & Frequência de Pedidos & Fração de ciclos em que a política emite pedido. \\
ADI & \textit{Average Demand Interval} & Intervalo médio entre períodos com demanda positiva. \\
CV$^2$ & Quadrado do coeficiente de variação & Medida de variabilidade relativa da demanda. \\
\textit{Lumpy} & Demanda grumosa & Regime com muitos períodos sem venda e alta variabilidade nos períodos com venda. \\
\textit{Smooth} & Demanda regular & Regime com alta frequência e baixa variabilidade. \\
\textit{Erratic} & Demanda errática & Regime com alta frequência e alta variabilidade. \\
\textit{Intermittent} & Demanda intermitente & Regime com muitos períodos sem venda e magnitude relativamente estável quando há venda. \\
RL & \textit{Reinforcement Learning} & Família de métodos usada para aprender políticas de reposição por interação com o ambiente simulado. \\
GA & \textit{Genetic Algorithm} & Meta-heurística usada para parametrizar políticas de reposição e compor políticas híbridas. \\
SA & \textit{Simulated Annealing} & Meta-heurística usada para otimização de parâmetros de políticas de reposição. \\
PSO & \textit{Particle Swarm Optimization} & Meta-heurística baseada em enxame usada para parametrizar políticas de reposição. \\
DE & \textit{Differential Evolution} & Meta-heurística evolutiva usada para otimização de parâmetros de políticas de reposição. \\
DQN & \textit{Deep Q-Network} & Agente usado para aprender uma política de reposição no ambiente de simulação. \\
PPO & \textit{Proximal Policy Optimization} & Agente usado para aprender uma política de reposição no ambiente de simulação. \\
SARSA & \textit{State-Action-Reward-State-Action} & Agente usado para aprender uma política de reposição no ambiente de simulação. \\
GA-DQN & Política híbrida GA + DQN & Política especialista que combina busca evolutiva e refinamento por aprendizado por reforço. \\
GA-PPO & Política híbrida GA + PPO & Política especialista que combina busca evolutiva e PPO. \\
EOQ & \textit{Economic Order Quantity} & Política clássica baseada em quantidade econômica de pedido. \\
$(s,S)$ & Política de ponto de pedido e nível máximo & Política que emite pedido quando o estoque cai abaixo de $s$, repondo até $S$. \\
Jornaleiro / \textit{Newsvendor} & Política de período único baseada em quantil da demanda & Política clássica usada como \textit{baseline} de reposição. \\
MASE & \textit{Mean Absolute Scaled Error} & Métrica de erro de previsão usada como sinal auxiliar. \\
sMAPE & \textit{Symmetric Mean Absolute Percentage Error} & Métrica secundária de erro de previsão. \\
RMSSE & \textit{Root Mean Squared Scaled Error} & Métrica escalada de erro quadrático de previsão. \\
\textit{Walk-forward} & Validação temporal progressiva & Protocolo que avalia modelos usando apenas dados passados para prever ciclos futuros. \\
Política dominante & Melhor política viável para uma série ou perfil & Política com menor CTI entre aquelas que atendem ao limiar mínimo de NS. \\
\textit{Policy label} & Rótulo de política & Classe supervisionada usada para treinar o PSE. \\
ASP & \textit{Algorithm Selection Problem} & Formulação teórica em que se escolhe o melhor algoritmo para cada instância do problema. \\
NFL & \textit{No Free Lunch} & Teorema que justifica a inexistência de uma política universalmente mais adequada para todos os perfis de demanda. \\
MEIO & \textit{Multi-Echelon Inventory Optimization} & Otimização de estoques em redes com múltiplos níveis ou escalões. \\
GSM & \textit{Guaranteed Service Model} & Família de modelos para posicionamento de estoques de segurança em redes multi-escalão. \\
\textit{Base-stock} & Política de estoque-base & Política que repõe o estoque até um nível-alvo predefinido. \\
SKU & \textit{Stock Keeping Unit} & Unidade individual de produto usada para controle operacional e análise de vendas. \\
MAPE & \textit{Mean Absolute Percentage Error} & Métrica percentual de erro de previsão, problemática quando há demanda zero. \\
M4/M5 & Competições internacionais de previsão & Benchmarks amplamente usados para comparar métodos de previsão de séries temporais. \\
ARIMA/SARIMAX & Modelos estatísticos de séries temporais & Modelos de previsão usados como alternativas auxiliares na etapa preditiva. \\
LSTM & \textit{Long Short-Term Memory} & Rede neural recorrente usada para capturar padrões sequenciais em séries temporais. \\
ANN & \textit{Artificial Neural Network} & Rede neural artificial usada como modelo de previsão de demanda. \\
XGBoost & \textit{Extreme Gradient Boosting} & Modelo baseado em árvores e gradiente \textit{boosting} usado para previsão e classificação. \\
LightGBM & \textit{Light Gradient Boosting Machine} & Modelo alternativo de \textit{boosting} considerado para classificação supervisionada. \\
RandomForest & Floresta aleatória & Modelo de ensemble baseado em múltiplas árvores de decisão. \\
DRL & \textit{Deep Reinforcement Learning} & Aprendizado por reforço profundo, isto é, RL com aproximação por redes neurais. \\
HAPPO & \textit{Heterogeneous-Agent Proximal Policy Optimization} & Variante multiagente de PPO discutida em trabalhos relacionados. \\
ELA & \textit{Exploratory Landscape Analysis} & Caracterização de instâncias por atributos de paisagem em seleção automática de algoritmos. \\
MILP/ILP & Programação inteira linear mista / inteira & Formulações exatas de otimização com variáveis inteiras e restrições lineares. \\
SB3 & \textit{Stable-Baselines3} & Biblioteca de referência para implementação de agentes de aprendizado por reforço. \\
API & \textit{Application Programming Interface} & Interface de programação usada para padronizar chamadas de software. \\
YAML & \textit{YAML Ain't Markup Language} & Formato de arquivo usado para parametrização dos experimentos. \\
Parquet & Formato colunar de dados & Formato eficiente para armazenar bases tabulares particionadas. \\
Kedro & \textit{Framework} de engenharia de dados & Ferramenta usada para organizar \textit{pipelines} reprodutíveis de processamento e simulação. \\
\textit{Dataset} & Conjunto de dados & Base organizada usada como entrada para análises, simulações ou treinamento. \\
\textit{Benchmark} & Protocolo comparativo & Avaliação padronizada usada para comparar políticas candidatas sob as mesmas condições. \\
\textit{Baseline} & Método de referência & Política ou modelo usado como ponto de comparação para avaliar ganhos relativos. \\
\textit{Feature} & Característica explicativa & Variável calculada a partir da série e usada como entrada de modelos de seleção ou previsão. \\
\textit{Pipeline} & Fluxo reprodutível de processamento & Sequência organizada de etapas de dados, simulação, avaliação e geração de artefatos. \\
\textit{Rolling horizon} & Horizonte deslizante & Protocolo temporal em que a janela de avaliação avança ao longo do tempo. \\
\textit{Rolling-origin} & Origem deslizante & Validação temporal em que o ponto de origem do treino avança progressivamente. \\
\textit{Burstiness} & Concentração de demanda em picos & Medida de irregularidade associada a eventos raros e intensos de demanda. \\
\textit{Bullwhip} & Efeito chicote & Amplificação da variabilidade de pedidos ao longo da cadeia de suprimentos. \\
Fronteira de Pareto & Conjunto de soluções não dominadas & Representa alternativas em que não é possível melhorar uma métrica sem piorar outra. \\
Teste de Wilcoxon & Teste não paramétrico pareado & Avalia diferenças entre políticas quando a normalidade não pode ser assumida. \\
Teste de Friedman & Teste não paramétrico para múltiplos métodos & Compara várias políticas simultaneamente por postos. \\
Teste de Nemenyi & Teste pós-hoc não paramétrico & Identifica quais pares diferem após rejeição no teste de Friedman. \\
Cohen's $d$ & Tamanho de efeito padronizado & Mede a magnitude prática da diferença entre duas políticas. \\
\bottomrule
\end{longtable}
\endgroup

\newpage
\lhead{\rightmark}
